\section{嵌套模型检验与变量筛选}
\srcnote{课件：第一讲第 20–24 页\quad 课堂代码：筛选自变量.R、stepwise.R、回归诊断.R 前半\quad 脚本：\texttt{scripts/03\_假设检验与变量筛选.R}\quad 数据：datasetC（n = 200，5 个候选变量）、diabetes27\_hw（n = 27）}

\subsection{增量 F 检验：anova(小模型, 大模型)}
所有“去掉某些变量行不行”的问题，都是比较两个嵌套模型。R 的通用写法是 \Verb|anova(小, 大)|：

\begin{RCode}[R · 偏 F 检验]
dd   <- read.csv("data/diabetes27_hw.csv")[, -1]
full <- lm(glucose ~ tc + tg + insulin + hba1c, data = dd)
m3   <- lm(glucose ~ tg + insulin + hba1c, data = dd)
anova(m3, full)                                       # 检验 tc
F_tc <- (deviance(m3) - deviance(full)) / 1 / (deviance(full) / df.residual(full))
pf(F_tc, 1, df.residual(full), lower.tail = FALSE)    # 手算，同上
coef(summary(full))[, "t value"]^2                    # 单个变量：偏 F = t²
anova(lm(glucose ~ insulin + hba1c, data = dd), full) # 同时检验 tc 与 tg
\end{RCode}
\begin{ROut}
> anova(m3, full)
Model 1: glucose ~ tg + insulin + hba1c
Model 2: glucose ~ tc + tg + insulin + hba1c
  Res.Df    RSS Df Sum of Sq      F Pr(>F)
1     23 89.454
2     22 88.841  1   0.61285 0.1518 0.7006
\end{ROut}
输出里 RSS 一列就是两个模型的残差平方和，$F=(89.454-88.841)/1\div(88.841/22)=0.152$。作业一要求“写出检验过程，包括 F 值、自由度、P 值”，三样都在这张表里：$F(1,22)=0.152$，$P=0.70$。

\subsection{drop1() 与 add1()}
\Verb|drop1(fit, test = "F")| 对模型里每个变量各做一次“删掉它”的偏 F 检验，\Verb|add1(fit, scope, test = "F")| 对候选变量各做一次“加上它”的检验。后退法的每一步看 \Verb|drop1|，前进法的每一步看 \Verb|add1|。不写 \Verb|test = "F"| 时只给 AIC，不给 P 值（课堂代码 \Verb|stepwise.R| 后半就是这种情况，所以要手算 F）。

\subsection{I 型、II 型、III 型平方和}
\begin{RCode}[R · 顺序与平方和类型]
anova(lm(glucose ~ tc + tg, data = dd))     # I 型：tc 的 SS 是它单独的贡献
anova(lm(glucose ~ tg + tc, data = dd))     # 顺序一换，两行都变
cor(dd$tc, dd$tg)                           # 0.63：两者相关才会这样
library(car)
Anova(lm(glucose ~ tc + tg, data = dd), type = "II")   # 与顺序无关
\end{RCode}
\begin{table}[htbp]
\centering\small
\caption{三种平方和的区别与 R 写法}
\begin{tabular}{@{}p{2.2cm}p{7.2cm}p{4.6cm}@{}}
\toprule
类型 & 每一行检验什么 & R\\
\midrule
I 型（顺序） & 在排在它前面的项之后新增的解释量 & \Verb|anova(fit)|\\
II 型 & 在其他所有\textbf{同级或更低级}项之后的贡献（不调整含它的交互项） & \Verb|car::Anova(fit, type = "II")|\\
III 型 & 在其他所有项（含交互）之后的贡献 & 先设和为零编码，再 \Verb|Anova(fit, type = "III")|\\
\bottomrule
\end{tabular}
\TableNote{无交互项时，II 型与 III 型相同，都等于 \texttt{drop1()} 的偏 F。平衡的析因设计中三者相同；不平衡时三者不同（第 6.9 节）。}
\end{table}

\subsection{为什么不能用 R² 选模型}
\begin{RCode}[R · 一组嵌套模型的五个指标]
dc <- read.csv("data/datasetC.csv")[, -1]
forms <- list("1" = glu ~ 1, "age" = glu ~ age, "age+bmi" = glu ~ age + bmi,
              "age+bmi+exercise" = glu ~ age + bmi + exercise,
              "+sleep" = glu ~ age + bmi + exercise + sleep,
              "full" = glu ~ age + bmi + exercise + sleep + waist)
fits <- lapply(forms, lm, data = dc)
t(sapply(fits, function(m) c(R2 = summary(m)$r.squared,
  adjR2 = summary(m)$adj.r.squared, AIC = AIC(m), BIC = BIC(m))))
\end{RCode}
\begin{ROut}
> t(sapply(fits, ...))
                     R2  adjR2      AIC      BIC
1                0.0000 0.0000 428.6045 435.2011
age              0.2791 0.2755 365.1542 375.0492
age+bmi          0.4191 0.4132 323.9729 337.1662
age+bmi+exercise 0.4692 0.4611 307.9146 324.4062
+sleep           0.4702 0.4594 309.5350 329.3249
full             0.4712 0.4576 311.1739 334.2621
\end{ROut}
$R^2$ 一路上升；校正 $R^2$、AIC、BIC 都在 age+bmi+exercise 处最优，与课件第 23 页的表一致。\Verb|lapply(列表, lm, data = dc)| 一次拟合一组模型，是比较多个模型时很省事的写法。

\subsection{全子集：leaps::regsubsets()}
\begin{RCode}[R · 全子集回归]
library(leaps)
all_sub <- regsubsets(glu ~ ., data = dc, nbest = 1)   # 每个变量个数下保留最好的 1 个
ss <- summary(all_sub)
ss$which                                     # 每行入选了哪些变量（TRUE/FALSE 矩阵）
cbind(adjR2 = ss$adjr2, Cp = ss$cp, BIC = ss$bic)
which.max(ss$adjr2); which.min(ss$cp); which.min(ss$bic)   # 都是第 3 行
coef(all_sub, 3)                             # 第 3 行模型的系数
plot(all_sub, scale = "adjr2")               # 黑格 = 入选
\end{RCode}
注意 \Verb|summary(regsubsets(...))| 返回的行号是模型的序号，不是变量个数（\Verb|nbest = 5| 时第 11 行是“3 个变量里排第 1”的模型）；课件笔记例 15-1 里 \Verb|which.max(...$adjr2)| 返回 11 就是这个意思。$C_p$ 的理想值是 $p+1$，本例 3 个变量时 $C_p=2.72$。\Verb|regsubsets| 返回的不是 lm 对象，不能直接 \Verb|vif()|、\Verb|summary()| 出系数表；选好后要用选中的变量重新 \Verb|lm()| 一次。

\subsection{按 AIC 的逐步法：step()}
\begin{RCode}[R · step 的四种用法]
m0    <- lm(glu ~ 1, data = dc)
mfull <- lm(glu ~ age + bmi + waist + sleep + exercise, data = dc)
step(mfull, direction = "backward")                                  # 后退
step(m0, scope = list(lower = ~ 1, upper = formula(mfull)), direction = "forward")
step(m0, scope = list(lower = ~ 1, upper = formula(mfull)), direction = "both")
step(m0, scope = list(lower = ~ 1, upper = formula(mfull)), direction = "both",
     k = log(nrow(dc)))                                              # 按 BIC
step(mfull, scope = list(lower = ~ age, upper = formula(mfull)))     # 强制保留 age
\end{RCode}
\begin{ROut}
> step(m0, scope = ..., direction = "forward")
Start:  AIC=-140.97
glu ~ 1
           Df Sum of Sq    RSS     AIC
+ age       1    27.310 70.542 -204.42
+ bmi       1    22.137 75.715 -190.27
+ waist     1    16.882 80.970 -176.85
+ exercise  1    10.647 87.206 -162.01
<none>                  97.853 -140.97
+ sleep     1     0.034 97.819 -139.04
\end{ROut}
读法：每一步列出“做这个动作后的 AIC”，\Verb|<none>| 是不动，取 AIC 最小的那一行，直到 \Verb|<none>| 最小为止。四种方式在本例都停在 age+bmi+exercise。几点要点：

\begin{itemize}
\item \Verb|scope| 的 \Verb|lower| 写进去的变量不会被删，混杂因素应当这样强制保留，而不是交给 AIC 决定。
\item \Verb|k = log(n)| 把惩罚从 2 换成 $\ln n$，就是按 BIC 选。
\item 这里的 AIC 是 \Verb|extractAIC()|。课堂代码 \Verb|回归诊断.R| 里那几行 \Verb|-140.97 + 200*(1+log(2*pi)) + 2| 就是把它换算成 \Verb|AIC()| 的数值：$-140.97+200(1+\ln2\pi)+2=428.60$，正是上表空模型的 AIC。
\end{itemize}

\subsection{按 P 值逐步：手工一步步走}
课件和作业用的是按 P 值的前进、后退、逐步法（$\alpha_{\text{入}}$、$\alpha_{\text{出}}$），R 的 \Verb|step()| 不做这个，可以用 \Verb|add1/drop1| 手工走，或者用 olsrr 自动走。

\begin{RCode}[R · 前进法与后退法（diabetes27）]
sc <- ~ tc + tg + insulin + hba1c
s0 <- lm(glucose ~ 1, data = dd)
add1(s0, sc, test = "F")                    # 第 1 步：hba1c 的 P 最小，进入
s1 <- update(s0, . ~ . + hba1c)
add1(s1, sc, test = "F")                    # 第 2 步：tc 进入（P = 0.03）
s2 <- update(s1, . ~ . + tc)
drop1(s2, test = "F")                       # 逐步法：进完一步要回头看能否删
add1(s2, sc, test = "F")                    # 第 3 步：最小 P = 0.088 > 0.05，停
drop1(full, test = "F")                     # 后退法：tc 的 P 最大（0.70），删
drop1(update(full, . ~ . - tc), test = "F") # 余下三个都 < 0.05，停
\end{RCode}
前进法（$\alpha=0.05$）停在 tc + hba1c，后退法停在 tg + insulin + hba1c——起点不同，终点就不同。后退法的循环可以写成一个函数，脚本 §3.7 的 \Verb|backward_p()| 就是作业一 bootstrap 用的那个。

\begin{RCode}[R · olsrr 自动逐步]
library(olsrr)
ols_step_forward_p(full, p_val = 0.05)         # 前进：p_val
ols_step_backward_p(full, p_val = 0.05)        # 后退：p_val
k <- ols_step_both_p(full, p_enter = 0.10, p_remove = 0.15)
k$model                                     # 最终模型（lm 对象，可以接着 summary）
\end{RCode}
olsrr（本机 0.7.1，课堂代码的写法也是如此）里前进、后退两个函数的阈值参数叫 \Verb|p_val|，逐步法才是 \Verb|p_enter|、\Verb|p_remove|。参数名写错不会报错，只会被 \Verb|...| 吞掉、悄悄用默认值 0.3，所以跑完要看一眼输出里的入选标准。

\subsection{筛选结果稳不稳：bootstrap}
\begin{RCode}[R · 重抽样 500 次看后退法的结果]
set.seed(2026)
sel <- replicate(500, {
  i <- sample(nrow(dd), replace = TRUE)
  paste(backward_p(dd[i, ], c("tc", "tg", "insulin", "hba1c"), "glucose"), collapse = "+")
})
head(sort(table(sel), decreasing = TRUE) / 500)
sapply(c("tc", "tg", "insulin", "hba1c"), function(v) mean(grepl(v, sel)))
\end{RCode}
\begin{ROut}
> head(sort(table(sel), decreasing = TRUE) / 500)
   tg+insulin+hba1c          tg+insulin            tc+hba1c               hba1c
              0.286               0.162               0.128               0.120
> sapply(...)           # 各变量被选入的比例
     tc      tg insulin   hba1c
  0.262   0.544   0.652   0.712
\end{ROut}
全样本选出的模型在重抽样中只出现 29\%。27 例、4 个相关的自变量，筛选结果本身就很不稳定，这正是课件说“筛过的 P 值要打折”的原因。\Verb|replicate(B, {...})| 加 \Verb|sample(n, replace = TRUE)| 是 bootstrap 最简单的写法。

\begin{notebox}[筛选的几条规矩]
先按专业知识定下必须调整的变量（放进 \Verb|scope| 的 \Verb|lower|），再让准则处理“可加可不加”的变量；写清起点、方向和准则；筛选后报告的 P 值和 CI 偏乐观，要在方法部分交代“变量经逐步法筛选”；样本小、自变量相关时，用 bootstrap 看一眼结果是否稳定。
\end{notebox}
